\documentclass[10pt,a4paper]{amsart}
\usepackage[margin=19mm]{geometry}
\usepackage{amsmath,amssymb,mathtools}
\usepackage[scr=rsfs,cal=euler]{mathalfa}
\usepackage{xfrac}
\usepackage[unicode,hidelinks,bookmarks=false]{hyperref}
\newcommand{\ie}{i.e.\ }
\newcommand{\cf}{cf.\ }
\newcommand{\resp}{respectively\ }
\newcommand{\Subsection}{\subsection}
\providecommand{\nobreakdash}{\nobreak\hbox{-}}
\setlength{\emergencystretch}{2em}
\multlinegap0pt
\usepackage[shortlabels]{enumitem}
\usepackage{mathtools}
\usepackage{xcolor}
\newtheorem{proposition}{Proposition}
\newtheorem{corollary}{Corollary}
\newtheorem{lemma}{Lemma}
\newtheorem{theorem}{Theorem}
\newcommand{\Ext}{\operatorname{Ext}}
\newcommand{\Gm}{\mathbf G_m}
\newcommand{\N}{\mathbf N}
\newcommand{\Ner}{\mathcal N}
\newcommand{\OK}{\mathcal O_K}
\newcommand{\Spec}{\operatorname{Spec}}
\newcommand{\Z}{\mathbf Z}
\newcommand{\cO}{\mathcal O}
\newcommand{\dd}{\mathfrak d}
\newcommand{\ord}{\operatorname{ord}}
\newcommand{\rad}{\operatorname{rad}}
\title{Ribet Points, geometric divisibility sequence and order of reductions on semiabelian varieties}
\author{Khai-Hoan Nguyen-Dang}
\date{Source: arXiv:2608.19742v2 (CC BY 4.0)}
\begin{document}
\maketitle

\noindent\emph{Source and licence.} Ribet Points, geometric divisibility sequence and order of reductions on semiabelian varieties. Version arXiv:2608.19742v2 (CC BY 4.0), distributed by arXiv under the Creative Commons Attribution 4.0 International licence, http://creativecommons.org/licenses/by/4.0/, as recorded in the arXiv licence metadata for this exact version and shown on the abstract page https://arxiv.org/abs/2608.19742v2. Full licence notice: \texttt{licenses/arxiv-2608.19742-CC-BY-4.0.txt}.\par\smallskip

\noindent\textit{Editorial scope.} Counted unit: the article's theorem thm:main (source0.tex lines 1478--1513). The article's complete original proof of it is delivered with it (source0.tex lines 1515--1576, one \texttt{proof} environment of 62 lines). No mathematics is written, repaired or re-derived: the statement and the proof are the article's own lines, in the article's own notation, and no retained line is substituted or re-flowed.  Retained uncounted as the source-local context the proof needs: lines 300--308; lines 842--873; lines 1166--1179; lines 1200--1220; lines 1290--1336; lines 1368--1377.  Nothing was omitted from any retained span, so the package carries no omission record.  No retained line contains a citation, so the wrapper carries no bibliography and no bibliography entry.  No retained line contains a figure environment, an image-inclusion command or drawing code, so no image is delivered and the package declares no asset dependency.  Editorial formatting only, in addition: an AMS \texttt{amsart} preamble that restates the article's own declarations and macros used by the retained lines, namely the environments \texttt{proposition}, \texttt{corollary}, \texttt{lemma}, \texttt{theorem} on separate counters, together with the standard packages those lines need.  The retained environments print in this document as Proposition 1, Corollary 1, Lemma 1, Theorem 1 in order of appearance, whereas the article prints the same items with the numbering of its own full document.  The complete native source of the article as distributed in the arXiv e-print for this exact version is bundled unchanged as \texttt{sources/208121/source0.tex}.
\begin{proposition}
\label{prop:intro-return-excludes-order}
Let $G/K$ be semiabelian, let $\Ner/\OK$ be its N\'eron lft-model, and let
$P\in G(K)$ have infinite order.  If $n>1$ and
\[
 \mathfrak d_{\Ner}(nP)=\mathfrak d_{\Ner}(P),
\]
then no finite place of semiabelian reduction satisfies $d_v(P)=n$.
\end{proposition}

\begin{proposition}
\label{prop:finite-cone-return}
Let $P\in G(K)$ have infinite order.  Suppose that there are a finite set $S$
of finite places and integers $h_1,\ldots,h_r>1$ such that $G$ has finite-type
semiabelian reduction outside $S$ and
\[
 d_v(P)\in h_1\N\cup\cdots\cup h_r\N
 \qquad(v\notin S).
\]
For $v\in S$, let
\[
 o_v:=\ord(\overline P_v)\in\N\cup\{\infty\},
 \qquad p_v:=\operatorname{char}(k_v),
\]
where the order is taken in the special fiber of the N\'eron lft-model.  Define
\[
 Q:=\rad\!\left(
 h_1\cdots h_r
 \prod_{\substack{v\in S\\1<o_v<\infty}}o_v
 \prod_{\substack{v\in S\\o_v=1}}p_v
 \right).
\]
Then $Q$ is squarefree and
\begin{equation}\label{eq:finite-cone-core}
 (n,Q)=1
 \quad\Longrightarrow\quad
 \dd_{\Ner}(nP)=\dd_{\Ner}(P).
\end{equation}
Consequently $\mathscr T(G,P)$ has lower natural density at least
$\varphi(Q)/Q$; every rational prime $s\nmid Q$ and every $a\ge1$ is a return
index.
\end{proposition}

\begin{corollary}\label{cor:ribet-reduction}
After deleting finitely many places of a number field $K$, all the data above
spread out and the normalized section commutes with reduction.  If
$m_v=\ord(\overline q_v)$, then
\[
 [m_v]\overline{R_\beta(q)}_v
 =e_{m_v}^A(\beta\overline q_v,\overline q_v),
 \qquad
 \ord(\overline{R_\beta(q)}_v)\mid m_v^2.
\]
These are scheme-theoretic identities on the finite flat group scheme
$A^\vee[m_v]$; therefore no prime-to-residue-characteristic hypothesis is
required.
\end{corollary}

\begin{lemma}\label{lem:amplification}
Let
\[
 0\longrightarrow T\longrightarrow H\xrightarrow{\pi}B\longrightarrow0
\]
be an exact sequence of finite abelian groups.  Let
\[
 R\in H,\qquad Q=\pi(R),\qquad t\in T,\qquad P=R+t.
\]
Fix a rational prime $\ell$, put $a=v_\ell(\ord(t))$, and suppose that
\begin{equation}\label{eq:order-control}
 v_\ell(\ord(R))
 \le c\,v_\ell(\ord(Q))+b
\end{equation}
for integers $c\ge1$ and $b\ge0$.  Then
\begin{equation}\label{eq:amplification-conclusion}
 v_\ell(\ord(P))
 \ge
 \max\left\{0,\left\lceil\frac{a-b}{c}\right\rceil\right\}.
\end{equation}
\end{lemma}

\begin{proposition}
\label{thm:abstract-shift}
Let
\[
 1\longrightarrow T\longrightarrow G\xrightarrow{\pi}B\longrightarrow0
\]
be an exact sequence of semiabelian varieties over a number field $K$, with
$T$ a torus.  Let $R\in G(K)$, put $Q_0=\pi(R)$, let $t\in T(K)$ be torsion,
and set $P=R+t$.  Assume that $P$ has infinite order.  Suppose that there is a
finite set $S$ such that, for every rational prime $\ell$, one has chosen
integers $c_\ell\ge1$ and $b_\ell\ge0$ satisfying
\begin{equation}\label{eq:abstract-local-control}
 v_\ell\bigl(\ord(\overline R_v)\bigr)
 \le c_\ell v_\ell\bigl(\ord(\overline Q_{0,v})\bigr)+b_\ell
 \qquad(v\notin S).
\end{equation}
Assume also that $t$ retains its generic order in every fiber outside $S$, and
define
\begin{equation}\label{eq:abstract-N}
 N_{\mathbf c,\mathbf b}(t):=\prod_\ell
 \ell^{\max\{0,\lceil(v_\ell(\ord t)-b_\ell)/c_\ell\rceil\}}.
\end{equation}
Only primes dividing $\ord(t)$ contribute.  Then, after enlarging $S$ if
necessary,
\[
 N_{\mathbf c,\mathbf b}(t)\mid d_v(P)
 \qquad(v\notin S).
\]
If $N_{\mathbf c,\mathbf b}(t)>1$, then:
\begin{enumerate}[label=\textup{(\roman*)}]
\item for every $n\ge1$ with $N_{\mathbf c,\mathbf b}(t)\nmid n$,
\[
 \dd_{\Ner}(nP)\cO_{K,S}=\cO_{K,S};
\]
\item there is a squarefree integer $Q$, divisible by
$\rad(N_{\mathbf c,\mathbf b}(t))$, such that
\begin{equation}\label{eq:abstract-coprime-return}
 (n,Q)=1\quad\Longrightarrow\quad
 \dd_{\Ner}(nP)=\dd_{\Ner}(P);
\end{equation}
\item $\mathscr T(G,P)$ has lower natural density at least $\varphi(Q)/Q$, and
for every rational prime $r\nmid Q$ and every $a\ge1$,
\[
 \dd_{\Ner}(r^{a}P)=\dd_{\Ner}(P).
\]
\end{enumerate}
\end{proposition}

\begin{lemma}
\label{lem:dense-lift}
Let
\[
 1\longrightarrow\Gm\longrightarrow G_q\xrightarrow{\pi}A\longrightarrow0
\]
be an extension over a field $K$ of characteristic zero, represented by a
non-torsion point $q\in A^\vee(K)$.  If $P\in G_q(K)$ has Zariski-dense
projection $\pi(P)$ in $A$, then $\mathbf ZP$ is Zariski dense in $G_q$.
\end{lemma}

\begin{theorem}
\label{thm:main}
Assume that $\mathbf Z(\delta q)$ is Zariski dense in $A$, and let
$\Ner/\OK$ be the N\'eron lft-model of $G_q$.  Then $G_q$ is geometrically
nonsplit, $P$ is non-torsion, and $\mathbf ZP$ is Zariski dense in $G_q$.
These geometric conclusions do not require $N_{\delta,t}>1$.

Assume in addition that $N_{\delta,t}>1$.  Then:
\begin{enumerate}[label=\textup{(\roman*)}]
\item there is a finite set $S$ of finite places such that
\begin{equation}\label{eq:N-divides-orders}
 N_{\delta,t}\mid d_v(P)\qquad(v\notin S);
\end{equation}
\item if $N_{\delta,t}\nmid n$, then
\begin{equation}\label{eq:localized-blanking}
 \dd_{\Ner}(nP)\cO_{K,S}=\cO_{K,S};
\end{equation}
\item there is a squarefree integer $Q=Q(G_q,P,S)$ satisfying
$\rad(N_{\delta,t})\mid Q$ and
\begin{equation}\label{eq:global-coprime-return}
 (n,Q)=1
 \quad\Longrightarrow\quad
 \dd_{\Ner}(nP)=\dd_{\Ner}(P);
\end{equation}
consequently $\mathscr T(G_q,P)$ has lower natural density at least
$\varphi(Q)/Q$ and contains the progression $1+Q\Z_{\ge0}$;
\item for every rational prime $r\nmid Q$ and every $a\ge1$,
\begin{equation}\label{eq:prime-power-return}
 \dd_{\Ner}(r^{a}P)=\dd_{\Ner}(P),
 \qquad
 r^a\notin\mathscr R(G_q,P).
\end{equation}
In particular, all but finitely many rational primes are both full return
indices and missing exact reduction orders.
\end{enumerate}
\end{theorem}

\begin{proof}
The density assumption and $\dim A>0$ imply that $q$ is non-torsion.  Under
$\Ext_K^1(A,\Gm)\simeq A^\vee(K)$, the geometric extension class is
$q_{\overline K}$, which is non-torsion and hence nonzero.  Therefore $G_q$ is
geometrically nonsplit.  Since $\pi(P)=\delta q$ is dense,
Lemma~\ref{lem:dense-lift} proves the geometric assertions.

Enlarge a finite set $S$ and put $U=\Spec\cO_{K,S}$.  Choose abelian schemes
$\mathcal A_U^\vee/U$ and $\mathcal A_U/U$ extending $A^\vee$ and $A$, and a
semiabelian scheme $\mathcal G_U/U$ extending $G_q$, so that $\beta$,
$\widehat\beta$, $\delta$, $q$, the normalized Poincar\'e biextension, the
normalized Ribet section, and $t$ all spread out compatibly over $U$.  Require
also that the extended map
$\delta:\mathcal A_U^\vee\to\mathcal A_U$ be an isogeny whose kernel is killed
by $e_\delta$ and that $t$ retain exact order $h$ in every geometric fiber.

Fix $v\notin S$ and put
\[
 m_v:=\ord(\overline q_v),
 \qquad
 n_v:=\ord(\delta\overline q_v).
\]
Since $[n_v]\overline q_v\in\ker\delta$,
\begin{equation}\label{eq:m-divides-en}
 m_v\mid e_\delta n_v.
\end{equation}
By Corollary~\ref{cor:ribet-reduction},
\begin{equation}\label{eq:R-order-m2}
 \ord(\overline R_v)\mid m_v^2.
\end{equation}
Consequently, for every rational prime $\ell$,
\begin{equation}\label{eq:R-vell-bound}
 v_\ell\bigl(\ord(\overline R_v)\bigr)
 \le2v_\ell(n_v)+2v_\ell(e_\delta).
\end{equation}

Since $H^1(k_v,\Gm)=0$, taking $k_v$-points gives an exact sequence
\[
 0\longrightarrow k_v^\times
 \longrightarrow\mathcal G_U(k_v)
 \longrightarrow\mathcal A_U(k_v)
 \longrightarrow0.
\]
Apply Lemma~\ref{lem:amplification} to this sequence
with
\[
 Q_0=\delta\overline q_v,
 \quad R=\overline R_v,
 \quad t=\overline t_v,
 \quad c=2,
 \quad b=2v_\ell(e_\delta).
\]
It follows that $v_\ell(d_v(P))\ge\rho_\ell(\delta,t)$ for every $\ell$,
proving (i).

The remaining assertions follow from Proposition~\ref{thm:abstract-shift} with these
control data.  More explicitly, (ii) is the reduction-order criterion applied
to \eqref{eq:N-divides-orders}; (iii) is
Proposition~\ref{prop:finite-cone-return} with $h_1=N_{\delta,t}$; and the equality in
(iv) follows from (iii).  The exclusion from $\mathscr R(G_q,P)$ follows from
Proposition~\ref{prop:intro-return-excludes-order}.
\end{proof}

\end{document}
